\section{\texorpdfstring{$(1,1,2)$}{(1,1,2)} 型代数の次数 \texorpdfstring{$1$}{1} 中心拡大}
\label{sec:central}

代数 $C$ の\emph{次数 $1$ の正則中心拡大}とは、
$t\in Z(A)_1$ が非零因子であり、次数付き同型 $A/(t)\cong C$ を持つ代数 $A$ をいう。
ここで $Z(A)$ は $A$ の中心である。Hilbert 級数の変数には $s$ を用いる。

\subsection{中心拡大から Ore データへ}

まず、後で反復にも用いる中心拡大の性質をまとめる。

\begin{lemma}\label{lem:central-lift}
$D$ を連結局所有限次数付き代数、$u\in Z(D)_e$ を非零因子とし、
$e\ge1$、$E=D/(u)$ とおく。
$E$ が両側 noetherian な $d$ 次元 AS-regular代数で AS index $\ell$ を持つなら、
$D$ は両側 noetherian な $d+1$ 次元 AS-regular代数であり、
\[
  H_D(s)=\frac{H_E(s)}{1-s^e},\qquad
  \ASind(D)=\ell+e
\]
が成り立つ。さらに、$e=1$ で $E$ の最小生成次数が $(1^m,2)$ なら、
$D$ の最小生成次数は $(1^{m+1},2)$ である。
\end{lemma}

\begin{proof}
正則中心拡大の議論 \cite[Propositions~2.1 and~2.3]{CKS} を、
本稿の重み付き次数に合わせて示す。完全列
\[
  0\longrightarrow D(-e)\xrightarrow{\cdot u}D
    \longrightarrow E\longrightarrow0
\]
から Hilbert 級数の等式を得る。従って $E$ の多項式増大度から $D$ の
有限 GK 次元も従う。右 $D$-加群として $\pd_D E=1$ である。
長さ $d$ の $E$-自由分解で $k$ を分解すると、各項の $D$ 上の射影次元は
高々 $1$ なので、短完全列に分けた次元移動により $\pd_D k\le d+1$ となる。
上の完全列に $\underline{\Ext}_D^*(-,D)$ を適用すると、
\[
  \underline{\Ext}_D^q(E,D)=
  \begin{cases}
    E(e)&q=1,\\
    0&q\ne1.
  \end{cases}
\]
したがって環変更のスペクトル系列は $q=1$ の行に退化し、
\begin{equation}\label{eq:central-ext-shift}
  \underline{\Ext}_D^{i+1}(k,D)
  \cong\underline{\Ext}_E^i(k,E)(e)
\end{equation}
を与える。$i=d$ の右辺は $k(\ell+e)\ne0$ だから $\pd_D k=d+1$ であり、
連結次数付き代数に対する $\gldim D=\pd_D k$ により大域次元も $d+1$ である。
左加群についても同じ計算が成立し、左右の AS 条件と AS index $\ell+e$ を得る。

Noetherian 性も正の次数だけを用いて持ち上げられる。
減少 filtration $F^pD=u^pD$ に対して、中心性と正則性から
$\operatorname{gr}_uD\cong E[U]$ である。ここで $U$ は中心変数で、
filtration 次数 $1$、内部次数 $e$ を持つ。
同次右 ideal $J\subset D$ に誘導 filtration を入れると、
$\operatorname{gr}_uJ$ は noetherian 環 $E[U]$ の同次右 ideal である。
その有限個の二重同次生成元を $J$ へ持ち上げる。
内部次数 $N$ の元をこれらの元で近似し、その差の $u$-adic order を
順次増やすと、order は $N/e$ 以下なので補正は有限回で停止する。
従って $J$ はそれらの持ち上げで生成される。
任意の右 ideal に対しては内部次数による増加 filtration を用い、
同次 ideal の生成元を持ち上げて最高次数を減らせば有限生成性が従う。
左 ideal にも同じ議論が適用できるので、$D$ は両側 noetherian である。

$e=1$ の場合を考える。
$E$ の最小生成元を同次的に $D$ へ持ち上げ、それらに $u$ を加える。
これらが $D$ を生成することは次数に関する帰納法で従う。実際、
同次元の $E$ での像を生成元の式として持ち上げると、差は $uD$ に属し、
$u$ を除いた係数の次数は $1$ 小さい。また $D_1\to E_1$ の核は $ku$ である。
次数 $2$ の生成元の持ち上げが $D_1$ の積で書けるなら、その像も $E_1$ の積で書け、
$E$ における最小性に反する。よって指定された最小生成次数を得る。
\end{proof}

\begin{theorem}\label{thm:central-ore}
$C$ を両側 noetherian な $3$ 次元 AS-regular代数で、
最小生成次数が $(1,1,2)$ であるものとする。
$A/(t)\cong C$ を次数 $1$ の正則中心拡大とする。
このとき $B=\langle A_1\rangle$ は両側 noetherian なquadratic Koszul AS-regular代数で、
$\gldim B=3$ である。$A_2/(A_1A_1)$ の基底を持ち上げる元 $w\in A_2$ を選ぶと、
ある $\sigma\in\Aut_{\mathrm{gr}}(B)$ と次数 $2$ の $\sigma$ 微分 $\delta$ により
\begin{equation}\label{eq:central-ore}
 A\cong B[w;\sigma,\delta],\qquad
 t\in Z(B)_1,\qquad \sigma(t)=t,\qquad\delta(t)=0
\end{equation}
と書ける。さらに、$R=B/(t)$ は量子平面または Jordan 平面であり、
商への誘導写像を上線で表すと
\begin{equation}\label{eq:central-quotient-ore}
 C\cong R[w;\bar\sigma,\bar\delta].
\end{equation}
その Hilbert 級数は
\[
 H_B(s)=\frac1{(1-s)^3},\qquad
 H_C(s)=\frac1{(1-s)^2(1-s^2)},\qquad
 H_A(s)=\frac1{(1-s)^3(1-s^2)}
\]
であり、$\gldim A=4$、$\ASind(A)=5$ である。

逆に、$B$ を両側 noetherian な二次 $3$ 次元 AS-regular代数とし、
$t\in Z(B)_1$ を非零因子で、$R=B/(t)$ が $2$ 次元 AS-regularとなるものとする。
$\sigma\in\Aut_{\mathrm{gr}}(B)$ と次数 $2$ の $\sigma$ 微分 $\delta$ が
$\sigma(t)=t$、$\delta(t)=0$ を満たすなら、$B[w;\sigma,\delta]$ は
\eqref{eq:central-quotient-ore} で定まる $(1,1,2)$ 型代数の次数 $1$ の正則中心拡大である。
\end{theorem}

\begin{proof}
主定理を $A$ と $C$ の双方に適用し、Ore 基底によって
$B/(t)$ を $\langle C_1\rangle$ と同定する。これにより、中心元の条件と商の AS-regular性を得る。

命題~\ref{prop:index-three} により $\ASind(C)=4$ である。
補題~\ref{lem:central-lift} から、$A$ は両側 noetherian な $4$ 次元 AS-regular代数で、
最小生成次数は $(1,1,1,2)$、AS indexは $5$ である。
定理~\ref{thm:ore} を $n=3$ で適用すると、
$B=\langle A_1\rangle$ はquadratic Koszul $3$ 次元 AS-regular代数となり、
$A=B[w;\sigma,\delta]$ を得る。

$t\in A_1=B_1$ かつ $t\in Z(A)$ なので $t\in Z(B)$ である。
$wt=tw$ と Ore 関係を比較すると
\[
   (\sigma(t)-t)w+\delta(t)=0.
\]
$A=\bigoplus_{j\ge0}Bw^j$ における係数の一意性から、
$\sigma(t)=t$ および $\delta(t)=0$ が従う。
したがって $\sigma$ と $\delta$ は $(t)\subset B$ を保ち、$R=B/(t)$ に
$\bar\sigma\in\Aut_{\mathrm{gr}}(R)$ と次数 $2$ の $\bar\sigma$ 微分 $\bar\delta$ を誘導する。
中心性と Ore 基底から
\[
  tA=\bigoplus_{j\ge0}tBw^j,\qquad
  B\cap tA=tB,\qquad
  A/tA=\bigoplus_{j\ge0}(B/tB)w^j
\]
となる。商における積は誘導された Ore 関係で与えられるので、
\eqref{eq:central-quotient-ore} は実際に次数付き代数同型である。

一方、$C$ に定理~\ref{thm:ore} を $n=2$ で適用すると、
$\langle C_1\rangle$ は二次 $2$ 次元 AS-regular代数である。
$B$ の $C$ における像はちょうど $\langle C_1\rangle$ であり、その核は上で示した $tB$ だから、
\[
 R=B/(t)\cong\langle C_1\rangle.
\]
$k=\C$ 上の二次 $2$ 次元 AS-regular代数の関係式は、基底変換によって
\begin{equation}\label{eq:central-plane}
  R\cong k\langle x,y\rangle/(yx-qxy-\varepsilon x^2),
  \qquad
  (q,\varepsilon)\in
  \{(q,0)\mid q\in k^\times\}\cup\{(1,1)\}
\end{equation}
の形になる（\cite[Section~2, following Corollary~2.2]{RZ}）。
これはそれぞれ量子平面と Jordan 平面であり、$H_R(s)=(1-s)^{-2}$ である。
$t$ は $B$ 上でも非零因子であるから $H_B=H_R/(1-s)$、
さらに Ore 基底から $H_A=H_B/(1-s^2)$ および $H_C=H_R/(1-s^2)$ を得る。

逆向きでは、Ore 拡大 $A=B[w;\sigma,\delta]$ の左 $B$ 基底を用いる。
$\sigma(t)=t$ と $\delta(t)=0$ より $wt=tw$ なので、$t$ は $A$ の中心に属する。
$t\sum_jb_jw^j=0$ なら各 $tb_j=0$ であるから各 $b_j=0$ となり、
$t$ は $A$ 上でも非零因子である。商を計算する上の議論もそのまま適用できる。
自己同型を用いる次数付き Ore 拡大は、両側 noetherian 性と AS-regular性を保ち、
大域次元を $1$ 増やす
（\cite[Sections~1.2--1.3 and Theorem~3.3]{LWW}）。
これを $B$ と $R$ に適用すると、$A$ と $C=R[w;\bar\sigma,\bar\delta]$ は
それぞれ $4$ 次元と $3$ 次元の AS-regular代数となる。
有限 GK 次元は Ore 基底による上記の Hilbert 級数からも確認できる。
また $C_2=R_2\oplus kw$ なので $w\notin\langle C_1\rangle$ であり、
$C$ の最小生成次数は $(1,1,2)$ である。
\end{proof}

\subsection{次数 \texorpdfstring{$1$}{1} 部分の明示的表示}

\begin{proposition}\label{prop:central-base-presentation}
定理~\ref{thm:central-ore} に現れる $(B,t)$ は、$x,y\in B_1$ の選択により
\begin{equation}\label{eq:central-base-presentation}
 B\cong
 \frac{k\langle t,x,y\rangle}
 {\bigl(tx-xt,\ ty-yt,\ yx-qxy-\varepsilon x^2-t(ax+by)-ct^2\bigr)}
\end{equation}
と表示できる。ここで $\deg t=\deg x=\deg y=1$、$a,b,c\in k$ であり、
$(q,\varepsilon)$ は \eqref{eq:central-plane} の条件を満たす。
逆に、これらの任意のパラメータに対して、\eqref{eq:central-base-presentation} は
両側 noetherian な二次 $3$ 次元 AS-regular整域を定める。
$t$ は正則中心元であり、その商は \eqref{eq:central-plane} の $R$ である。
\end{proposition}

\begin{proof}
まず表示の右辺を $P$ とおく。可換多項式環 $S=k[t,x]$ 上で
\begin{align*}
  &\alpha(t)=t,\qquad &&\alpha(x)=qx+bt,\\
  &\beta(t)=0,\qquad &&\beta(x)=\varepsilon x^2+atx+ct^2
\end{align*}
と定める。$q\ne0$ なので $\alpha$ は次数付き自己同型であり、
$\alpha^{-1}(x)=q^{-1}(x-bt)$ である。
$\beta$ を自由代数から $\alpha$ 微分として延長すると、
$f=\varepsilon x^2+atx+ct^2$ に対して
\[
  \beta(xt-tx)=ft-tf=0\quad\text{in }S
\]
であるから、$S$ 上の $\alpha$ 微分を定める。その次数は $1$ である。
Ore 関係 $ys=\alpha(s)y+\beta(s)$ を $s=t,x$ について書くと
\eqref{eq:central-base-presentation} の関係式に一致するため、
\begin{equation}\label{eq:central-base-iterated-ore}
 P\cong k[t,x][y;\alpha,\beta].
\end{equation}
従って
\[
   \{t^i x^j y^h\mid i,j,h\ge0\}
\]
は $P$ の $k$ 基底であり、$H_P(s)=(1-s)^{-3}$ である。
さらに、$y$ の最高次の係数がそれぞれ $r_m,s_n\in S\setminus\{0\}$ である
非零元の積の最高次係数は $r_m\alpha^m(s_n)$ なので、$P$ は整域である。
$S$ は両側 noetherian な $2$ 次元 AS-regular代数であり、$\alpha$ は自己同型だから、
\cite[Sections~1.2--1.3 and Theorem~3.3]{LWW} により
$P$ は両側 noetherian な $3$ 次元 AS-regular代数である。
有限 GK 次元は既に計算した Hilbert 級数で確認される。
また、関係式から $t$ は中心元であり、PBW 基底からその乗法は単射である。
$t=0$ とすれば商は \eqref{eq:central-plane} に一致する。

次に、定理~\ref{thm:central-ore} の $B$ に戻る。
$R_1$ の基底を \eqref{eq:central-plane} の形に選び、$x,y\in B_1$ へ持ち上げる。
このとき $t,x,y$ は $B_1$ の基底である。
$yx-qxy-\varepsilon x^2$ は $B_2$ に属し、$R$ における像が $0$ なので、
ある $a,b,c\in k$ によって
\[
  yx-qxy-\varepsilon x^2=t(ax+by+ct)
\]
と書ける。したがって同じパラメータの $P$ から $B$ への次数付き全射がある。
両者の Hilbert 級数は $(1-s)^{-3}$ であるため、各次数で全射は同型となり、
$P\cong B$ が従う。
\end{proof}

\subsection{同型によるデータの同一視}

\eqref{eq:central-base-presentation} のパラメータと、その $B$ 上の
$(\sigma,\delta)$ は、生成元の選択に依存する。この依存は次の同値関係で記述できる。
中心直線を指定した拡大のデータを $(B,kt,\sigma,\delta)$ とし、
これに $A=B[w;\sigma,\delta]$ を対応させる。
別のデータを $(B',kt',\sigma',\delta')$、対応する代数を
$A'=B'[w';\sigma',\delta']$ とする。

\begin{proposition}\label{prop:central-isomorphisms}
中心直線を保つ次数付き同型 $F:A\to A'$、$F(kt)=kt'$ は、
次のデータと $1$ 対 $1$ に対応する：
次数付き同型 $\phi:B\to B'$、$\lambda,\mu\in k^\times$ および $h\in B'_2$ であって、
\begin{align}
  \phi(t)&=\mu t',\label{eq:central-iso-t}\\
  \phi\sigma&=\sigma'\phi,\label{eq:central-iso-sigma}\\
  \phi\delta(b)&=\lambda\delta'\phi(b)+h\phi(b)-\phi\sigma(b)h
       \qquad(b\in B)\label{eq:central-iso-delta}
\end{align}
を満たすものである。この対応における同型は
\[
  F|_B=\phi,\qquad F(w)=\lambda w'+h
\]
で与えられる。
\end{proposition}

\begin{proof}
命題~\ref{prop:ore-isomorphism} により、同型 $F$ の形と
\eqref{eq:central-iso-sigma}、\eqref{eq:central-iso-delta} の必要十分性を得る。
$t\in B_1$ なので、$F(kt)=kt'$ という条件はちょうど
\eqref{eq:central-iso-t} である。
\end{proof}

従って次数 $1$ の中心直線を指定した拡大の同型類は、
定理~\ref{thm:central-ore} のデータを
\eqref{eq:central-iso-t}--\eqref{eq:central-iso-delta} で同一視したものと一致する。
中心元そのものを指定する場合は $\mu=1$ とする。
固定した代数 $C$ への商写像まで指定する場合は、
$A/(t)\cong C$ という同型をデータに加え、誘導される商の同型がこの同一視と可換になることを課す。
中心直線を指定しない代数の同型を考える場合は、\eqref{eq:central-iso-t} を外せばよい。

\subsection{中心拡大の反復}

\begin{corollary}\label{cor:iterated-central}
$n\ge2$ とし、$C_0$ を定理~\ref{thm:central-ore} の $(1,1,2)$ 型代数とする。
連結局所有限次数付き代数の列 $C_0,C_1,\ldots,C_{n-2}$ があり、各 $j=1,\ldots,n-2$ について
\[
  t_j\in Z(C_j)_1\text{ は非零因子},\qquad
  C_j/(t_j)\cong C_{j-1}
\]
が成り立つとする。このとき $A=C_{n-2}$ は両側 noetherian な AS-regular代数で、
\[
 \gldim A=n+1,\qquad
 \ASind(A)=n+2,\qquad
 H_A(s)=\frac1{(1-s)^n(1-s^2)}
\]
を満たし、最小生成次数は $(1^n,2)$ である。従って定理~\ref{thm:ore} が適用され、
\[
 A=\langle A_1\rangle[w;\sigma,\delta],\qquad \deg w=2,
\]
と書ける。ここで $\langle A_1\rangle$ はquadratic Koszul $n$ 次元 AS-regular代数である。
\end{corollary}

\begin{proof}
$C_0$ の大域次元は $3$、AS indexは $4$、最小生成次数は $(1,1,2)$ であり、
定理~\ref{thm:central-ore} の証明で得た Ore 表示から
$H_{C_0}(s)=((1-s)^2(1-s^2))^{-1}$ である。
補題~\ref{lem:central-lift} を各段階に適用すると、
$j$ 段階目では大域次元は $3+j$、AS indexは $4+j$、
最小生成次数は $(1^{2+j},2)$ となり、Hilbert 級数には $(1-s)^{-j}$ が掛かる。
$j=n-2$ とすれば主張の数値と主定理のすべての仮定が得られる。
\end{proof}
